\section{Conclusion}
\label{sec:conclusion}

Across $43$ benchmarks, two years after Clay~v1.5's release, specialized 2025--2026 models lead: optical models (DINOv3) on pure-RGB benchmarks, and multi-temporal Sentinel models on phenology benchmarks. Clay~v1.5 leads two $0.1\text{--}1\,\text{m}$ segmentation tasks under a frozen linear probe. It reaches $91\text{--}94\%$ of the best EO model on an unseen $21$-band instrument and a $98\%$ median of the best model on an independent leaderboard's $15$ classification datasets. It is a permissive Apache-2.0 encoder for sub-meter, multi-band, or unseen-instrument inputs without retraining. We draw three design recommendations for next-generation sensor-agnostic EO models. (1)~Pair dynamic wavelength stems with explicit cross-sensor and vision--language alignment for open-vocabulary retrieval and descriptive tasks, and model multi-date sequences natively rather than reducing time semantics to post-hoc derivatives of single-snapshot embeddings. (2)~Decouple attention stride from dense readout resolution via sub-patch poolers~\cite{olmoearth2026v13}. (3)~Retain Matryoshka supervision through the final distillation stage.

{\small
\paragraph{Data, Code, and Acknowledgments.}
\label{sec:acknowledgments}
Apache-2.0 code (\path{Clay-foundation/model}, \path{Clay-foundation/stacchip}), Apache-2.0 weights (\path{huggingface.co/made-with-clay/Clay}), CC-BY-4.0 embedding catalogs (\path{source.coop/clay}), and companion evaluation scripts (\path{Clay-foundation/clay-v1.5-paper}) are publicly available. \textbf{Funding and competing interests.} Clay is a program of Renaissance Philanthropy. Clay v0 through Clay v1.5 were developed with support from Schmidt Futures, the Grantham Foundation for the Protection of the Environment, and Amazon Web Services (AWS) while Clay was a fiscally sponsored project of Radiant Earth (501(c)(3)). During that period, Ode contributed to model development, strategy, operations, and the development of applications and demos. Development Seed contributed to building and training the Clay models. Legend AI funded its authors' work on this report, including the evaluation. AWS, the Schmidt Family Foundation, the Monterey Bay Aquarium, and Amazon Sustainability also funded applied projects. Legend AI builds commercial products on Earth embeddings, including those produced from Clay models. M.G., D.H., A.C., and L.C. are affiliated with Ode, which does applied work with Earth embeddings. We thank \textbf{Wei Ji Leong}, \textbf{Lilly Thomas}, \textbf{Leo Thomas}, \textbf{Ryan Avery}, \textbf{Myscon Truong}, \textbf{Mansour Raad}, \textbf{Jeff Albrecht}, and \textbf{Nat Manning}. We thank \textbf{Romeo Kienzler} for his contributions to packaging Clay as a \texttt{pip}-installable Python library. We also thank open-source contributors Aimee Barciauskas, Dan Bonomo, Kevin Booth, Isaac Campbell, Michele Claus, Isaac Corley, Chuck Daniels, Beno\^it de Chateauvieux, Tyler Erickson, Aninda Goswamy, Riddhi Goswami, Ash Hoover, Kelsey Josund, Freddie Kalaitzis, Maxime Lenormand, Javier Marin Tur, Bill Morris, Beatriz Peres, Ra\'ul Ramos, Ferdinand Schenck, Peter Sherman, Jed Sundwall, Jonas Viehweger, Lauren Yee, and Brayden Zhang, and the TorchGeo and TerraTorch maintainers (as of November 2026).

\section*{Responsible AI Use Disclosure}
\label{sec:ai_disclosure}
In November 2026 the authors used Gemini Carbon, Claude Opus, Claude Fable, and ChatGPT Astra for research, drafting, and adversarial review. The authors take full responsibility for this work.
}
